% Source: https://artofproblemsolving.com/community/c3h2445348_help_needed_for_compound_interest_problem
% Stable identity: aops-topic-2445348-q01
Let $r=1.085$ and let $P$ be the payment. After the first year Stefan owes $12000r-P$; after the second year he owes $(12000r-P)r-P=12000r^2-Pr-P$; continuing in the same way, after the fourth payment he owes
\[12000r^4-P\left(r^3+r^2+r+1\right).\]
This must be $0$. By the formula for a geometric series, $r^3+r^2+r+1=\dfrac{r^4-1}{r-1}$, so
\[P=\dfrac{12000r^4(r-1)}{r^4-1}=\dfrac{1020r^4}{r^4-1}.\]
Now $r^2=1.177225$ and $r^4=1.385858700625$, so
\[P=\dfrac{1020\cdot1.385858700625}{0.385858700625}=\dfrac{1413.5758746375}{0.385858700625}\approx3663.4547.\]
Each payment is $\boxed{\$3663.45}$.

(Check, with amounts rounded to the cent: $13020.00-3663.45=9356.55$; $10151.86-3663.45=6488.41$; $7039.92-3663.45=3376.47$; $3663.47-3663.45=0.02$, which is only the rounding error from the cents.)
